% =============================================================================
% CHAPTER 4 — AI and NLP Foundations (TalentOs)
% =============================================================================
% Include from main thesis with (typical):
%   \usepackage{amsmath,amssymb,amsfonts}
%   \usepackage{graphicx}
%   \usepackage{hyperref}
%   \usepackage{minitoc}   % only if you use \minitoc
% Place figure as: figures/ia.png  OR set \graphicspath{{figures/}}
% =============================================================================

\chapter{AI and NLP Foundations}
\label{chap:ai_nlp_foundations}

\minitoc
\newpage

% -----------------------------------------------------------------------------
\section*{Introduction}
\addcontentsline{toc}{section}{Introduction}
% -----------------------------------------------------------------------------

This chapter establishes the theoretical and technical groundwork required to understand the \textbf{TalentOs} intelligence core. By exploring the hierarchy of Artificial Intelligence (AI) and the mechanisms of Natural Language Processing (NLP), we define how unstructured human narratives are transformed into structured, actionable data for modern recruitment workflows.

% -----------------------------------------------------------------------------
\section{Hierarchical Evolution of AI}
\label{sec:ch4:hierarchical_evolution}
% -----------------------------------------------------------------------------

The modernization of recruitment processes relies on a technological ecosystem designed to bridge the gap between human professional experience and machine-readable data. In the context of \textbf{TalentOs}, this transformation is governed by a nested structure of technologies that work in synergy to convert unstructured \textbf{CVs and job descriptions} into actionable hiring insights.

The evolution of these systems has moved toward a \textbf{Cognitive Recruitment} approach. This paradigm emphasizes a \textbf{human-in-the-loop} model, where machine-supported reasoning assists recruiters in navigating complex data without replacing the essential human judgment required for final hiring decisions.

\subsection{Artificial Intelligence, Machine Learning, Deep Learning, and Large Language Models}
\label{sec:ch4:ai_ml_dl_llm}

To establish the technical framework of \textbf{TalentOs}, it is essential to define the layers of intelligence as they are actually implemented within the platform.

\begin{figure}[htbp]
    \centering
    \includegraphics[width=0.8\textwidth]{figures/ia.png}% adjust path if needed
    \caption{Hierarchical relationship between AI, ML, DL, and LLMs.}
    \label{fig:ch4:ai_hierarchy}
\end{figure}

\begin{itemize}
    \item \textbf{Artificial Intelligence (AI):} This serves as the overarching ambition to build systems capable of perception, reasoning, and action under uncertainty. In \textbf{TalentOs}, AI provides the global logic for automating the recruitment pipeline—from parsing candidate data to supporting multi-actor workflows—enabling the system to handle scale and complexity that would be overwhelming for manual processing alone.

    \item \textbf{Machine Learning (ML):} A data-driven paradigm where models are induced from examples. Within the platform, ML acts as the mathematical basis for ranking and similarity: rather than relying solely on rigid keyword rules, TalentOs utilizes \textbf{transformer-based dense embeddings} and \textbf{cosine similarity} between requirement and CV representations (e.g., skill phrases). Geometrically, for embedding vectors $\mathbf{u}, \mathbf{v} \in \mathbb{R}^{d}$,
    \begin{equation}
        \mathrm{sim}_{\cos}(\mathbf{u}, \mathbf{v}) \;=\;
        \frac{\mathbf{u}^{\top}\mathbf{v}}{\lVert \mathbf{u}\rVert_{2}\,\lVert \mathbf{v}\rVert_{2}}
        \;=\;
        \frac{\sum_{i=1}^{d} u_i v_i}{\sqrt{\sum_{i=1}^{d} u_i^{2}}\;\sqrt{\sum_{i=1}^{d} v_i^{2}}}\,.
        \label{eq:ch4:cosine}
    \end{equation}
    A similarity score is not, by itself, a hiring decision: it becomes actionable when combined with thresholds, weighting across skill/education/experience/language signals, and human validation (see \autoref{sec:ch4:evaluation}).

    \item \textbf{Deep Learning (DL):} A specialized subset of ML based on multi-layered neural networks. In TalentOs, \textbf{deep learning underpins the transformer encoders} used to produce dense text embeddings (e.g., via \textbf{SentenceTransformer}-class models) and the \textbf{large language models} used for extraction and assistance. In other words, DL supplies the representation-learning capacity that makes semantic matching and contextual parsing effective.

    \item \textbf{Large Language Models (LLMs):} Representing the current frontier of generative and instruction-following behaviour within the DL family, LLMs serve as the contextual core for several TalentOs services. The platform leverages hosted inference (e.g., \textbf{Groq}) to perform \textbf{information extraction}, mapping professional narratives into complex \textbf{JSON} schemas while exploiting contextual cues that are difficult to capture with brittle rules alone.
\end{itemize}

\subsection{Natural Language Processing: Bridging the Communication Gap}
\label{sec:ch4:nlp_bridge}

Natural Language Processing (NLP) serves as the critical interface between human professional narratives and machine logic. It addresses the central difficulty of recruitment: workflows are mediated primarily by \textbf{language} and, in later stages, by \textbf{multimodal} interaction (e.g., interview audio/video alongside transcripts).

In \textbf{TalentOs}, NLP transforms linguistic signals into structured representations—such as \textbf{tokens}, \textbf{semantic embeddings}, and \textbf{cross-section context}. This process supports:

\begin{enumerate}
    \item \textbf{Semantic interpretation:} Moving beyond literal keyword matching so that professional intent is preserved (e.g., relating \textit{DevOps} to \textit{Site Reliability Engineering}).

    \item \textbf{Structural standardization:} A pipeline that distinguishes \textbf{digital PDF text} from \textbf{scanned pages}; the system applies \textbf{OCR} when extracted text is insufficient, followed by \textbf{LLM-assisted structuring} into validated fields.

    \item \textbf{Knowledge augmentation (RAG):} Integrating retrieval with generation so recruiters can query the candidate pool in natural language while answers are \textbf{anchored} to retrieved evidence linked to job requirements.
\end{enumerate}

% -----------------------------------------------------------------------------
\section{Foundational Mechanisms of Modern NLP}
\label{sec:ch4:foundational_nlp}
% -----------------------------------------------------------------------------

\subsection{Tokenization and Embeddings: Digitizing and Vectorizing Text}
\label{sec:ch4:tokenization_embeddings}

Raw text is not directly usable as a fixed-length numerical input for most modern models; learning algorithms operate on discrete symbols within controlled vocabularies. Two processes bridge this gap:

\begin{itemize}
    \item \textbf{Tokenization:} Splitting text into \textit{tokens} (commonly subword units). Subword algorithms mitigate out-of-vocabulary issues and trade off vocabulary size, maximum sequence length, and multilingual behaviour.

    \item \textbf{Embeddings:} Mapping tokens or entire passages to vectors in $\mathbb{R}^{d}$ so that semantic proximity corresponds (approximately) to geometric proximity. Cosine similarity is a standard measure; it is the same quantity formalized in \autoref{eq:ch4:cosine}.
\end{itemize}

\paragraph{TalentOs link.}
Candidate and job texts are embedded using \textbf{transformer-based encoders}. The reference configuration uses a \textbf{SentenceTransformer} model (\textit{all-MiniLM-L6-v2}) producing \textbf{384}-dimensional vectors. These representations power semantic skill matching and support vector indexing for retrieval-augmented workflows.

\subsection{Transformer Architecture and Attention}
\label{sec:ch4:transformer_attention}

Transformers address limitations of earlier sequence models for long-range dependencies and parallel training. The core mechanism is \textbf{scaled dot-product self-attention}. Given query, key, and value matrices $Q$, $K$, and $V$, with key dimension $d_k$,
\begin{equation}
    \mathrm{Attention}(Q,K,V)
    \;=\;
    \mathrm{softmax}\!\left(\frac{QK^{\top}}{\sqrt{d_k}}\right)V\,.
    \label{eq:ch4:attention}
\end{equation}
In practice, models stack \textbf{multi-head attention}, \textbf{residual connections}, \textbf{layer normalization}, and \textbf{position-wise feed-forward} blocks. \textbf{Encoder-style} models are strong for embeddings and retrieval; \textbf{decoder-style} LLMs support autoregressive generation and schema-constrained outputs under prompting.

\paragraph{TalentOs link.}
Encoder-style embeddings support matching and dense retrieval, while decoder LLMs (via hosted APIs such as \textbf{Groq}) support CV-to-JSON extraction and recruiter-assistant style generation, typically with conservative decoding settings for factual tasks.

\subsection{Large Language Models: Capabilities and Limits}
\label{sec:ch4:llm_capabilities_limits}

\subsubsection*{Capabilities}

\begin{itemize}
    \item \textbf{Instruction following and open-ended language behaviour} for both understanding and generation.
    \item \textbf{In-context adaptation} via prompts and examples without updating base weights at deployment time.
    \item \textbf{Structured outputs} (e.g., JSON modes) enabling reliable consumption by backend services.
\end{itemize}

\subsubsection*{Limits and mitigations}

\begin{table}[htbp]
\centering
\begin{tabular}{|p{3.2cm}|p{10.8cm}|}
\hline
\textbf{Topic} & \textbf{Brief description and mitigation in TalentOs} \\
\hline
Hallucinations &
Plausible but unsupported statements. Mitigated using \textbf{RAG grounding}, explicit ``answer from context'' instructions, schema validation, and human review for critical decisions. \\
\hline
Knowledge cut-off &
Static pretraining corpora. Mitigated by retrieving \textbf{fresh private data} (CVs, jobs, applications) at inference time. \\
\hline
Context window &
Hard token limits. Mitigated by \textbf{chunking}, summarization where appropriate, and retrieval of only relevant passages. \\
\hline
Latency and cost &
Managed via hosted inference and engineering choices (batching, caching where applicable); empirical latency should be reported with measured percentiles when evaluating the system. \\
\hline
\end{tabular}
\caption{LLM limitations and corresponding mitigation strategies in TalentOs.}
\label{tab:ch4:llm_limits}
\end{table}

\paragraph{TalentOs link.}
LLMs are used for schema-oriented extraction and assistant-style interactions. \textbf{RAG} and \textbf{vector stores} (e.g., \textbf{Chroma}) ground outputs in retrieved evidence and reduce reliance on stale parametric knowledge alone.

% -----------------------------------------------------------------------------
\section{Automated Document Processing and Information Extraction}
\label{sec:ch4:document_processing}
% -----------------------------------------------------------------------------

Recruitment data rarely arrives as a clean relational record: curricula vitae are distributed as PDFs and images with heterogeneous layouts. The extraction pipeline therefore separates \textbf{document acquisition and text recovery} from \textbf{semantic structuring}, so downstream matching and decision-support operate on consistent representations.

\subsection{OCR: Recovery from Legacy CV Formats}
\label{sec:ch4:ocr}

\begin{itemize}
    \item \textbf{Why OCR matters:} Scanned CVs, design-heavy exports, and low-quality captures may lack reliable selectable text.
    \item \textbf{PDF reality:} Differentiate \textit{digital PDFs} (embedded text) from \textit{image-like pages}. TalentOs follows a \textbf{native text extraction first} strategy and applies OCR as a \textbf{fallback} when extracted text is too short or unreliable.
    \item \textbf{OCR process and limitations:} OCR outputs plain text but may lose complex layout (multi-column tables), and accuracy depends on resolution, fonts, skew, and noise.
\end{itemize}

\paragraph{TalentOs link.}
The platform combines fast native text extraction with an OCR fallback on rendered page images when needed, so candidates are not excluded solely by document format.

\subsection{CV Parsing and Schema-Guided Extraction}
\label{sec:ch4:cv_parsing}

\begin{itemize}
    \item \textbf{Conceptual framing:} Classical NER assigns token-level labels (e.g., person/organization spans). In TalentOs, extraction is best described as \textbf{schema-guided parsing}: mapping free text into a validated \textbf{JSON schema} (skills, experience, education, contact, certifications, etc.).
    \item \textbf{Goals:} Enable semantic matching, retrieval chunking, UI population, and traceability (linking extracted claims to source text where possible).
    \item \textbf{Methodology:} Beyond brittle regular expressions, TalentOs uses \textbf{LLM-based structured extraction} with contextual disambiguation (e.g., distinguishing a degree mention from colloquial uses of ``master'').
\end{itemize}

\paragraph{Mitigations and validation.}
JSON response constraints reduce format failures; field-level checks catch common inconsistencies; critical hiring decisions remain subject to human review.

\paragraph{Terminology note.}
Throughout this document, ``NER'' refers to \textbf{named-entity-style fields} recovered through \textbf{schema-based LLM extraction}, rather than a standalone sequence-labelling NER model trained on an annotated token corpus—unless explicitly stated otherwise.

% -----------------------------------------------------------------------------
\section{Augmented Intelligence and Deployment Strategies}
\label{sec:ch4:augmented_intelligence}
% -----------------------------------------------------------------------------

While preceding sections explain how text is recovered, embedded, and structured, this section explains how TalentOs augments generative components with \textbf{private, up-to-date evidence} through retrieval and vector search, and how behaviour is shaped at inference time through prompting. Parameter-efficient fine-tuning techniques (e.g., LoRA/QLoRA) are discussed as \textbf{future extensions} beyond the current implementation focus.

\subsection{Retrieval-Augmented Generation and Private Recruitment Corpora}
\label{sec:ch4:rag}

LLM deployment in recruitment must address static weights, limited context, hallucination risk, and privacy. \textbf{RAG} retrieves relevant passages at query time, injects them into the prompt, and conditions generation on that evidence.

\begin{itemize}
    \item \textbf{Chunking:} Documents are split (e.g., recursive splitting with controlled \texttt{chunk\_size} and overlap) to stabilize retrieval granularity.
    \item \textbf{Indexing:} Chunks are embedded in the same semantic space used for matching (\autoref{eq:ch4:cosine}).
    \item \textbf{Grounding:} Prompts instruct the model to rely on retrieved context and to acknowledge missing evidence.
\end{itemize}

\paragraph{TalentOs link.}
The RAG service combines \textbf{LangChain} components (\textbf{ChatGroq}, text splitters) with a \textbf{Chroma} vector store scoped to recruitment artefacts (e.g., per-job corpora), enabling recruiter questions to be answered with retrieved candidate/job context.

\subsection{Vector Indexes, Dense Retrieval, and Semantic Matching}
\label{sec:ch4:vector_indexes}

Vector indexes serve two related functions:

\begin{itemize}
    \item \textbf{Dense retrieval (RAG):} retrieve top-$k$ chunks most similar to a query embedding.
    \item \textbf{Candidate--job matching:} compare embeddings of requirements and CV fields (skills, etc.) using cosine similarity signals combined with additional category logic.
\end{itemize}

\paragraph{TalentOs link.}
\textbf{Chroma} stores embedded chunks with metadata for traceability; \textbf{SentenceTransformer}-class embeddings align retrieval and matching semantics beyond keyword overlap.

\subsection{Parameter-Efficient Fine-Tuning: LoRA and QLoRA (Future Work)}
\label{sec:ch4:lora_qlora}

\textbf{Low-Rank Adaptation (LoRA)} trains small adapter matrices while freezing most base weights; \textbf{QLoRA} reduces memory via quantized base weights. These methods are \textbf{not required} to explain the current TalentOs deployment, which emphasises retrieval, structured extraction, and prompt discipline—but they are relevant \textbf{future work} for domain-specialised behaviour on curated HR corpora with controlled governance.

\subsection{Prompt Design and Output Control}
\label{sec:ch4:prompt_engineering}

Prompting is the primary interface for shaping frozen-model behaviour:

\begin{itemize}
    \item \textbf{Role instructions} (system/user messages) defining parser vs.\ assistant roles.
    \item \textbf{Schema constraints} (JSON modes) for reliable backend consumption.
    \item \textbf{RAG prompts} that restrict answers to retrieved context and use conservative decoding (e.g., low temperature) for factual tasks.
\end{itemize}

% -----------------------------------------------------------------------------
\section{Evaluation Metrics and Validation Protocol}
\label{sec:ch4:evaluation}
% -----------------------------------------------------------------------------

Quality assurance cannot rely on a single scalar. \textbf{Cosine similarity} quantifies alignment in embedding space; \textbf{precision, recall, and F1} apply only after a module is formalised as a prediction task with explicit ground truth—whether for retrieved chunks, candidate--job suitability labels, or extracted schema fields. This section defines metrics and a validation protocol; \textbf{if experimental results are produced}, they should be reported with dataset size, labelling protocol, and uncertainty (confidence intervals or inter-rater agreement where applicable).

\subsection{Objectives and Scope of Quality Assurance}
\label{sec:ch4:qa_scope}

For TalentOs, quality is interpreted along complementary axes:

\begin{itemize}
    \item \textbf{Correctness and robustness} of extraction under OCR noise and CV diversity.
    \item \textbf{Grounding} of assistant outputs to retrieved evidence.
    \item \textbf{Fairness considerations} in matching design (e.g., emphasising competence signals rather than sensitive attributes—aligned with product policy).
    \item \textbf{Operational quality}: latency percentiles, error rates, and JSON validity rates measured on real or held-out samples.
\end{itemize}

\subsection{Metrics for Semantic Matching and Retrieval}
\label{sec:ch4:metrics_matching_retrieval}

\paragraph{Cosine similarity as a signal.}
Similarity from \autoref{eq:ch4:cosine} can be converted into a binary decision with a threshold $\tau$ when evaluating against labelled pairs $(\mathbf{u},\mathbf{v},y)$ with $y\in\{0,1\}$:
\begin{equation}
    \hat{y} \;=\;
    \begin{cases}
        1, & \mathrm{sim}_{\cos}(\mathbf{u}, \mathbf{v}) \ge \tau,\\[0.35em]
        0, & \mathrm{otherwise}.
    \end{cases}
    \label{eq:ch4:cosine_threshold}
\end{equation}
Thresholds may be chosen by maximising accuracy or $F_1$ on a validation set, subject to fairness review.

\paragraph{Ranking metrics for retrieval (RAG).}
Let $r_k$ be the number of relevant items among the top-$k$ retrieved, and $R_{\mathrm{tot}}$ the total relevant items in the collection. Then
\begin{align}
    \mathrm{Precision@}k &= \frac{r_k}{k}, \\
    \mathrm{Recall@}k &= \frac{r_k}{R_{\mathrm{tot}}}.
\end{align}
For a query $i$, let $\mathrm{rank}_i$ be the rank of the first relevant retrieved item (if none exists, define the reciprocal contribution as $0$). Over $Q$ queries,
\begin{equation}
    \mathrm{MRR}
    \;=\;
    \frac{1}{Q}\sum_{i=1}^{Q}\frac{1}{\mathrm{rank}_i}\,.
    \label{eq:ch4:mrr}
\end{equation}

\subsection{Metrics for Structured Extraction (JSON)}
\label{sec:ch4:metrics_extraction}

For a binary prediction task with confusion counts $TP$, $FP$, and $FN$ for the positive class, precision, recall, and the $F_1$-score are
\begin{equation}
    P \;=\; \frac{TP}{TP+FP},
    \qquad
    R \;=\; \frac{TP}{TP+FN},
    \qquad
    F_1 \;=\; \frac{2PR}{P+R}\,.
    \label{eq:ch4:prf1}
\end{equation}

Let fields be indexed by $f=1,\dots,F$, with confusion counts $(TP_f, FP_f, FN_f)$ defined per field against a gold standard. Micro-averaging aggregates counts before computing precision and recall:
\begin{align}
    P_{\mathrm{micro}} &= \frac{\sum_{f} TP_f}{\sum_{f} (TP_f + FP_f)}, \\
    R_{\mathrm{micro}} &= \frac{\sum_{f} TP_f}{\sum_{f} (TP_f + FN_f)}, \\
    F_{1,\mathrm{micro}} &= \frac{2 P_{\mathrm{micro}} R_{\mathrm{micro}}}{P_{\mathrm{micro}}+R_{\mathrm{micro}}}\,.
\end{align}
For strict identifiers (email/phone), \textbf{exact match rate} is often reported separately. \textbf{JSON validity rate} measures syntactic conformance to the enforced schema output channel.

\subsection{LLM Output Quality and Groundedness}
\label{sec:ch4:metrics_llm}

Beyond field extraction, assistant outputs benefit from:

\begin{itemize}
    \item \textbf{Groundedness checks:} verifying claims against retrieved chunks (supported/unsupported).
    \item \textbf{Hallucination audits:} manual or rubric-based review on sampled queries.
    \item \textbf{Cut-off stress tests:} queries requiring post-training facts, where RAG should supply evidence from the database rather than parametric memory.
\end{itemize}

\subsection{Human-in-the-Loop Validation}
\label{sec:ch4:human_validation}

Final recruitment decisions remain human-governed. Validation may include:

\begin{itemize}
    \item agreement between AI-ranked shortlists and recruiter rankings on sampled cases,
    \item an error taxonomy (OCR failure, chunking boundary errors, semantic drift, formatting),
    \item controlled prompt comparisons (A/B) against a fixed evaluation set.
\end{itemize}

% End of chapter ----------------------------------------------------------------
